\colorlet{darkgreen}{green!50!black}
% \usebeamercolor{structure}
% %\colorlet{color1}{beamer@structure.fg}
% \definecolor{color1}{named}{fg}
% \usebeamercolor{normal text}
\definecolor{color1}{HTML}{3333B3}
\colorlet{color2}{red!80!black}
\colorlet{color3}{green!40!black}
\colorlet{myblue}{color1}
\colorlet{myorange}{color2}
\colorlet{mygreen}{color3}
%\colorlet{myorange}{orange!80!black}
\colorlet{myred}{red!65!black}

\colorlet{myblue1}{myblue!15!white}
\colorlet{myblue2}{myblue!7!white}
\colorlet{mygreen1}{mygreen!15!white}
\colorlet{myorange1}{myorange!15!white}

\colorlet{nodegreen}{green!70!black!20}
\colorlet{importantnodecolor}{orange!10}

\colorlet{semialgebraiccolor}{mygreen}

\tikzset{reset node/.style={text width={},minimum height=0pt,minimum width=0pt, anchor=center, line width=1pt, text color=normal text.fg}}

\tikzset{
  myarrow/.style={-latex',thick}
}

\tikzset{
  plotsoa/.style={thick,color=myblue,mark=x},
  plotnew1/.style={thick,color=myred,mark=+},
}

\tikzset{
  plotqh/.style={thick,color=myred,mark=+},
  ploth/.style={thick,color=myblue,mark=x},
  finalpoint/.style={draw,fill,inner sep=1.5pt}
}


\tikzset{
  polysystem/.style={draw,rectangle,
    % rounded corners,fill=white
    color=structure.fg,thick,%
    % minimum height=0.8cm,%
    align=center,anchor=center}
}

\colorlet{darkgreen}{green!50!black}

% \tikzset{
%   darkgreen/.style={draw=green!50!black}
% }
\tikzset{
  covered/.style={color=structure.fg!20}
}
\tikzset{
  algorithm/.style={draw=darkgreen, text=darkgreen,rectangle,rounded corners=5pt,outer sep=3pt,thick,fill=white}
}
\tikzset{
  complexity/.style={draw,myred,rectangle,rounded corners,outer sep=3pt,thick,fill=white}
}

\newcommand\pgfdrawmonomialgrid[3]{% %% xmax, ymax, style
  \def\borderoffset{0.5}
  \draw [thick, ->] (- \borderoffset,0) -- (#1 + \borderoffset,0) node[right] {$\deg_{X_{1}}$};
  \draw [thick, ->] (0, - \borderoffset) -- (0,#2 + \borderoffset) node[above] {$\deg_{X_{2}}$};
  \foreach \x in {0,1,...,#1}{
    \foreach \y in {0,1,...,#2}{
      \draw (\x,\y) node[#3] {};
    }
  }
}

\newcommand{\pgfdrawweightedmonomials}[5]{% %% xmax, ymax, weightx,
                                          % %% weighty, style
  \foreach \x in {0,#3,...,#1}{
    \foreach \y in {0,#4,...,#2}{
      \node[#5] at (\x,\y) {};
    }
  }
}

\newcommand{\pgfdrawdegreeline}[3]{% %% d, texte, style
  \draw[#3] (0,#1) -- ($(#1,0)+(0.5,-0.5)$) node[anchor=north] {#2};
  }

\tikzset{%
  mark empty monomial/.style={inner sep=2pt,outer sep=0pt,minimum size=3pt,
    circle, draw,myblue,transform shape}%
}
\tikzset{%
  mark full monomial/.style={mark empty monomial, myred, inner sep=2pt,
    fill}%
}
\tikzset{%
  nothing/.style={}}

\tikzset{
  invisible/.style={opacity=0},
  visible on/.style={alt=#1{}{invisible}},
  alt/.code args={<#1>#2#3}{%
    \alt<#1>{\pgfkeysalso{#2}}{\pgfkeysalso{#3}}
  },
  covered on/.style={alt=#1{opacity=0.2}{}}
}


\newcommand{\nodetrue}[2]{%
  \node [circle,draw,green,thick] at (#1,#2) {\checkmark}
}


\newcommand{\nodefalse}[2]{%
  \node [circle,draw,myred,thick] at (#1,#2) {X}
}





% For correct externalization with overlays

\makeatletter
\newcommand*{\overlaynumber}{\number\beamer@slideinframe}

\tikzset{
  beamer externalizing/.style={%
    execute at end picture={%
      \tikzifexternalizing{%
        \ifbeamer@anotherslide
        \pgfexternalstorecommand{\string\global\string\beamer@anotherslidetrue}%
        \fi
      }{}%
    }%
  },
  external/optimize=false
}

%\def\ifundefined#1{\expandafter\ifx\csname#1\endcsname\relax}

\let\orig@tikzsetnextfilename=\tikzsetnextfilename
\renewcommand\tikzsetnextfilename[1]{%
  % \ifundefined{tikzsetnextfilename}%
  % \fi%
  \orig@tikzsetnextfilename{#1-\overlaynumber}}
\makeatother

%\renewcommand{\tikzsetnextfilename}[1]{}

\tikzset{every picture/.style={beamer externalizing}}

\tikzset{external overlay/.style={external/export=false,overlay}}

% Catcodes in matrices
\tikzset{beamer matrix/.style={ampersand replacement=\&}}



\tikzset{plotnode/.style={fill=white,shape=circle,inner sep=0pt,font=\small}}
\tikzset{cntnode/.style={fill=white,shape=circle,inner sep=1pt,font=\small,draw}}


%%% Horizontal and vertical lines in tikz matrices
%%% See: http://tex.stackexchange.com/a/204437/9517

% \mvline[<style>]{<matrix name>}{<row number on the right hand side of the line>}
\newcommand\mvline[3][]{%
  \pgfmathtruncatemacro\hc{#3-1}
  \draw[#1]({$(#2-1-#3)!.5!(#2-1-\hc)$} |- #2.north) -- ({$(#2-1-#3)!.5!(#2-1-\hc)$} |- #2.south);
}
% \mhline[<style>]{<matrix name>}{<column number below of the line>}{<number of columns in a row>}
\newcommand\mhline[4][]{%
  \node[fit=(#2-#3-1),inner sep=0pt,outer sep=0pt](R){};
  \foreach \i in {1,...,#4}\node[fit=(R) (#2-#3-\i),inner sep=0pt,outer sep=0pt](R){};
  \draw[#1] (R.north -| #2.west) -- (R.north -| #2.east);
}


\tikzset{myarrow-decoration/.style 2 args={decoration={markings,mark=at position #1 with {\arrow[thick]
        {#2}}}}}

\tikzset{textnode/.style={align=left,text width=0.52\textwidth}}



\tikzset{
  plotbad/.style={very thick,color=myred%,mark=x
  },
  plotgood/.style={very thick,color=mygreen% ,mark=+
  },
  plotmaybe/.style={very thick,color=myorange% ,mark=o
  }
}

\tikzset{highlight/.style={%fill=yellow!70!white
    fill=green!15!white
    , rectangle, rounded corners},
  alerthighlight/.style={%fill=yellow!70!white
    fill=myorange!10!white
    , rectangle, rounded corners}}


\def\myleftmargin{2pt}
\def\mytopmargin{0.9cm}
\def\myrightmargin{2pt}
\def\mybottommargin{2pt}

\newcommand{\setcoordinatesoverlay}{
  \coordinate (page-nw) at ($(current page.north west) + (\myleftmargin,-\mytopmargin)$);
  \coordinate (page-n) at ($(current page.north) + (0,-\mytopmargin)$);
  \coordinate (page-ne) at ($(current page.north east) + (-\myrightmargin,-\mytopmargin)$);
  \coordinate (page-e) at ($(current page.east) + (-\myrightmargin,0)$);
  \coordinate (page-se) at ($(current page.south east) + (-\myrightmargin,\mybottommargin)$);
  \coordinate (page-s) at ($(current page.south) + (0,\mybottommargin)$);
  \coordinate (page-sw) at ($(current page.south west) + (\myleftmargin,\mybottommargin)$);
  \coordinate (page-w) at ($(current page.west) + (\myleftmargin,0)$);
  \path (page-nw) -- (page-se) node[pos=0.5] (page-center) {};
  \coordinate (page-c) at (page-center);
}

\newcommand{\setcoordinatesnooverlay}{
  \def\mywidth{128mm}
  \def\myheight{96mm}
  \coordinate (page-nw) at (\myleftmargin,\myheight-\mytopmargin);
  \coordinate (page-ne) at (\mywidth-\myrightmargin,\myheight-\mytopmargin);
  \coordinate (page-sw) at (\myleftmargin,\mybottommargin);
  \coordinate (page-se) at (\mywidth-\myrightmargin,\mybottommargin);
  \path[use as bounding box] (page-nw) -- (page-ne) node[pos=0.5,coordinate] (page-n) {}
  -- (page-se) node[pos=0.5,coordinate] (page-e) {}
  -- (page-sw) node[pos=0.5,coordinate] (page-s) {}
  -- (page-nw) node[pos=0.5,coordinate] (page-w) {}
  -- (page-se) node[pos=0.5,coordinate] (page-center) {};
  \coordinate (page-c) at (page-center);
}

\newcommand{\drawcross}{
  \begin{scope}[every path/.style={draw,line width=0.5pt,cyan!20!white},
    every node/.style={fill=white,inner sep=0pt}]
    \path (page-nw) node[anchor=north west] {NW} -- (page-se) node[anchor=south east] {SE};
    \path (page-ne) node[anchor=north east] {NE} -- (page-sw) node[anchor=south west] {SW};
    \path (page-w) -- (page-e);
    \path (page-n) -- (page-s);
    \node at (page-c) {C};
  \end{scope}
}  

\newcommand{\setcoordinates}{\setcoordinatesoverlay}


%%% Local Variables: 
%%% mode: latex
%%% TeX-master: "main"
%%% End: 
